\documentclass[12pt]{article}
\providecommand{\DeclareUnicodeCharacter}[2]{}% guard: bqf-paper uses it; XeTeX/tectonic may not define it
\PassOptionsToPackage{table}{xcolor}% bqf-paper loads xcolor; pre-pass the table option to avoid a clash
\usepackage{bqf-paper}
\usepackage{algorithm}
\usepackage{algpseudocode}
\usepackage{amsmath}
\usepackage[utf8]{inputenc} % allow utf-8 input
\usepackage[T1]{fontenc}    % use 8-bit T1 fonts
\usepackage{hyperref}       % hyperlinks
\usepackage{url}            % simple URL typesetting
\usepackage{booktabs}       % professional-quality tables
\usepackage{amsfonts}       % blackboard math symbols
\usepackage{nicefrac}       % compact symbols for 1/2, etc.
\usepackage{microtype}      % microtypography
\usepackage{xcolor}         % colors
\usepackage{tikz}
\usetikzlibrary{arrows.meta,positioning,calc,decorations.pathreplacing,shapes.geometric}
\usepackage[table]{xcolor}
\usepackage{listings}

% bqf-paper.sty has \endinput at line 308, so macros defined after it are not loaded.
\providecommand{\KG}{\mathcal{K}_G}
\providecommand{\rhostar}{\rho_*}
\providecommand{\Prob}{\mathbb{P}}
\providecommand{\sgn}{\operatorname{sgn}}
\providecommand{\arsinh}{\operatorname{arsinh}}
\providecommand{\erf}{\operatorname{erf}}
\newcommand{\He}{\mathsf{H}}                 % orthonormal probabilist Hermite
\newcommand{\Dop}{\mathrm{D}}              % Euler operator t d/dt (override sty's \Dop)
\newcommand{\Kr}{K}                          % bivariate Gaussian kernel
\newcommand{\gam}{\gamma}
\providecommand{\cB}{\mathcal B}

\title{\paperTitleSmallCaps{A degree--$3/5$ sign rounding that beats Krivine's bound}}
\author{}
\date{}

\begin{document}
\maketitle

\begin{abstract}
We construct an odd $\pm1$ rounding pair whose arcsine--normalized correlation
$H$ is almost linear, and we certify in interval arithmetic that the associated
Krivine admissibility radius exceeds that of the hyperplane.  This gives
\[
   \KG \;\le\; 1.7818666069360661
   \;<\; \frac{\pi}{2\log(1+\sqrt2)} \;=\; 1.7822139781913691,
\]
an improvement of at least $3.4737\cdot10^{-4}$ over Krivine's bound.  The proof
replaces the long inverse--coefficient estimate of earlier schemes by two short
ingredients: a Wiener--algebra contraction that reads the inverse majorant off
the correlation series itself, and a single scalar Sobolev inequality that
controls the coefficient tail.  The Sobolev norm is bounded through an exact
reduction of $H$ to a two--dimensional Gaussian integral and a rigorous
quadrature on the circle.  All inequalities are verified with \texttt{Arb}.
\end{abstract}

\section{Setup and main result}\label{sec:setup}

We use the normalizations of the surrounding paper.  For odd measurable
$f,g\colon\R^d\to\{-1,1\}$ and standard Gaussian vectors $(\bx,\by)$ in $\R^d$
with coordinatewise correlation $t\in(-1,1)$, write
\[
   \mathfrak H_{f,g}(t)=\E\bigl[f(\bx)\,g(\by)\bigr],
   \qquad
   H_{f,g}(t)=\frac{\pi}{2}\,\mathfrak H_{f,g}(t).
\]
The hyperplane rounding has $H_{\mathrm{hp}}(t)=\arcsin t$, with admissibility
radius
\[
   \rhostar=\arsinh(1)=\log(1+\sqrt2)=0.8813735870195430,
\]
and Krivine's bound is $\KG\le \pi/(2\rhostar)=1.7822139781913691$.

Throughout, $\He_n$ is the orthonormal probabilist Hermite polynomial,
$\E[\He_m(Z)\He_n(Z)]=\delta_{mn}$ for $Z\sim N(0,1)$; in particular
$\He_0=1$, $\He_1(x)=x$, and
\begin{equation}\label{eq:He3}
   \He_3(x)=\frac{x^3-3x}{\sqrt6}.
\end{equation}
If $(X,Y)$ are standard Gaussian with correlation $t$ then
$\E[\He_m(X)\He_n(Y)]=\delta_{mn}t^n$.

\subsection{The scheme}

Fix the three real numbers
\begin{equation}\label{eq:params}
   \eta=0.136419125,\qquad s_3=0.34101124,\qquad s_5=0.05276111,
\end{equation}
and set
\begin{equation}\label{eq:VW}
   V=1+s_3^2+s_5^2,\qquad \alpha=\frac{\eta}{\sqrt V}.
\end{equation}

\begin{definition}[Limiting correlation]\label{def:H}
Let $(Z,X,R_3,R_5)$ and $(Z',X',R_3',R_5')$ be jointly Gaussian, each coordinate
marginally $N(0,1)$, mutually independent except for
\[
   \mathrm{corr}(Z,Z')=\mathrm{corr}(X,X')=t,\qquad
   \mathrm{corr}(R_3,R_3')=t^3,\qquad \mathrm{corr}(R_5,R_5')=t^5 .
\]
Define
\begin{align}
   f &= \sgn\bigl(Z+\eta\,\He_3(X)+s_3R_3+s_5R_5\bigr),\label{eq:f}\\
   g &= \sgn\bigl(Z-\eta\,\He_3(X)-s_3R_3+s_5R_5\bigr),\label{eq:g}
\end{align}
and let $H(t)=\tfrac{\pi}{2}\,\E[fg']$, where $g'$ is $g$ in the primed
variables.  We write $H(t)=\sum_{m\ge1}b_mt^m$; only odd $m$ occur.
\end{definition}

The two sign patterns matter: the degree--$3$ term is anti--aligned (opposite
signs in $f$ and $g$) and the degree--$5$ term is aligned.  The correlations
$t,t^3,t^5$ encode the rule that $R_3$ and $R_5$ carry, respectively, third-- and
fifth--degree information; \cref{def:H} is the fixed--degree limit of a genuine
finite--dimensional construction in the sense of BMMN, and a strict limiting
certificate transfers to a finite scheme of dimension $d_N=2+2N$ for all large
$N$ (\cref{rem:transfer}).

\subsection{Main theorem}

\begin{theorem}\label{thm:main}
For the rounding pair of \cref{def:H} with parameters \eqref{eq:params},
\[
   \KG \;\le\; \frac{\pi}{2\gam}=1.7818666069360661,
   \qquad \gam=0.881545409 ,
\]
and hence $\KG$ is smaller than Krivine's bound by at least
$3.4737\cdot10^{-4}$.  Every inequality below is certified in interval
arithmetic.
\end{theorem}

\paragraph{Proof plan.}
The argument has four steps, each a short certified inequality.

\begin{enumerate}[leftmargin=2.2em]
\item \emph{Admissibility from the correlation series} (\cref{sec:wiener}).
A Wiener--algebra contraction shows that if the nonlinear mass of $H$ is small
enough, namely
\[
   \eta_0:=\sum_{m\ge3}\abs{b_m}\quad\text{and}\quad
   \eta_1:=\sum_{m\ge3}m\abs{b_m}
\]
satisfy $\gam+\eta_0<b_1$ and $\eta_1<b_1$, then $\KG\le\pi/(2\gam)$.  No inverse
series is needed.
\item \emph{The coefficients} (\cref{sec:reduction}).  An exact reduction turns
$H$ into a two--dimensional Gaussian integral and gives a finite, interval--checkable
formula for each $b_m$.  In particular $b_1>0.8815738220$, $b_3$ and $b_5$ are
of order $10^{-9}$ and $10^{-11}$, and the head sums of $\eta_0,\eta_1$ are
certified directly.
\item \emph{The tail} (\cref{sec:sobolev}).  A single scalar Sobolev norm
$B_3=\norm{\Dop^3H}_{L^2(\T)}$ controls $\sum_{m>N}\abs{b_m}$ through
Cauchy--Schwarz; the power weights $m^{-6}$ make the tail negligible.
\item \emph{Bounding $B_3$} (\cref{sec:b3}).  We bound $B_3$ by an averaged
Minkowski inequality applied to a combined differential kernel, reducing it to a
rigorous quadrature on the circle of certified Gaussian integrals.  This yields
$B_3\le 54.4641$, which is far more than enough.
\end{enumerate}

\noindent The numbers are assembled in \cref{sec:certificate}; the explicit
formulas behind every certified quantity are collected in
\cref{sec:formulas}.

\section{The admissibility criterion}\label{sec:wiener}

Let $\cA$ be the Wiener algebra of power series $U(w)=\sum_{n\ge0}u_nw^n$ with
$\norm{U}_{\cA}=\sum_n\abs{u_n}<\infty$; it is a Banach algebra,
$\norm{UV}_{\cA}\le\norm{U}_{\cA}\norm{V}_{\cA}$.  Write
\begin{equation}\label{eq:H-split}
   H(t)=b_1t+E(t),\qquad E(t)=\sum_{m\ge3}b_mt^m,
\end{equation}
and recall $\eta_0=\sum_{m\ge3}\abs{b_m}$, $\eta_1=\sum_{m\ge3}m\abs{b_m}$.

\begin{lemma}[Wiener contraction]\label{lem:wiener}
Suppose $b_1>0$ and
\begin{equation}\label{eq:wiener-cond}
   \eta_1<b_1,\qquad \gam+\eta_0<b_1
\end{equation}
for some $\gam>0$.  Then $H^{-1}$ is holomorphic on the disk of radius $\gam$ and
\[
   \sum_{n\ge1}\abs{A_n}\gam^n<1,\qquad H^{-1}(\zeta)=\sum_{n\ge1}A_n\zeta^n .
\]
Consequently $\KG\le\pi/(2\gam)$.
\end{lemma}

\begin{proof}
Put $U(w)=H^{-1}(\gam w)$.  The inverse equation $H(U)=\gam w$ is the fixed point
\[
   U=\Phi(U):=\frac{\gam w-E(U)}{b_1}.
\]
Let $\cB=\{U\in\cA:U(0)=0,\ \norm{U}_{\cA}\le1\}$.  If $U\in\cB$ then
$\norm{E(U)}_{\cA}\le\sum_{m\ge3}\abs{b_m}\norm{U}_{\cA}^m\le\eta_0$, so
$\norm{\Phi(U)}_{\cA}\le(\gam+\eta_0)/b_1<1$, i.e.\ $\Phi$ maps $\cB$ into
itself.  For $U,V\in\cB$, $\norm{U^m-V^m}_{\cA}\le m\norm{U-V}_{\cA}$, hence
$\norm{E(U)-E(V)}_{\cA}\le\eta_1\norm{U-V}_{\cA}$ and
$\norm{\Phi(U)-\Phi(V)}_{\cA}\le(\eta_1/b_1)\norm{U-V}_{\cA}$.  By
\eqref{eq:wiener-cond}, $\Phi$ is a contraction of $\cB$; its unique fixed
point is the scaled inverse, and
$\norm{U}_{\cA}=\sum_n\abs{A_n}\gam^n\le(\gam+\eta_0)/b_1<1$.  The bound
$\KG\le\pi/(2\gam)$ is Krivine's preprocessing criterion applied to a $\gam$ at
which the inverse majorant is below $1$.
\end{proof}

So \cref{thm:main} reduces to certifying $b_1$, $\eta_0$ and $\eta_1$.

\section{The correlation coefficients}\label{sec:reduction}

The four Gaussian directions of $f$ collapse to two.  Since $Z+s_3R_3+s_5R_5$ is
a single $N(0,V)$ variable independent of $X$, the function $f$ depends on its
arguments only through the pair
\[
   (W,X),\qquad W=\frac{Z+s_3R_3+s_5R_5}{\sqrt V}\sim N(0,1),
\]
and $f=\sgn\!\bigl(W+\alpha\,\He_3(X)\bigr)$ with $\alpha$ as in \eqref{eq:VW}.

\begin{definition}[Threshold and coefficient integrals]\label{def:Aab}
For $a,b\ge0$ let $F(w,x)=\sgn\!\bigl(w+\alpha\,\He_3(x)\bigr)$ and
\begin{equation}\label{eq:Aab}
   A_{a,b}=\E_{W,X}\bigl[\He_a(W)\He_b(X)\,F(W,X)\bigr].
\end{equation}
For $u\in\R$ set $q_a(u)=\E_R[\sgn(u+R)\He_a(R)]$, $R\sim N(0,1)$.
\end{definition}

\begin{lemma}[Closed form for $q_a$]\label{lem:qa}
With $\varphi(u)=(2\pi)^{-1/2}e^{-u^2/2}$,
\[
   q_0(u)=\erf\!\Bigl(\tfrac{u}{\sqrt2}\Bigr),\qquad
   q_a(u)=2\varphi(u)\,\frac{\He_{a-1}(-u)}{\sqrt a}\quad(a\ge1).
\]
Moreover $A_{a,b}=\E_X\bigl[\He_b(X)\,q_a(\alpha\,\He_3(X))\bigr]$, a
one--dimensional integral, and $A_{a,b}=0$ unless $a+b$ is odd.
\end{lemma}

\begin{proof}
The case $a=0$ is $\Prob[R>-u]-\Prob[R<-u]=\erf(u/\sqrt2)$.  For $a\ge1$,
$\E\He_a(R)=0$ gives $q_a(u)=2\int_{-u}^\infty\He_a(r)\varphi(r)\,dr$, and
$\frac{d}{dr}\bigl(\varphi(r)\He_{a-1}(r)\bigr)=-\sqrt a\,\varphi(r)\He_a(r)$
yields the stated form.  Integrating $W$ first in \eqref{eq:Aab} replaces
$\E_W[\He_a(W)\,\sgn(W+\alpha\He_3(X))]$ by $q_a(\alpha\He_3(X))$.  Under
$X\mapsto-X$ one has $\He_3(-x)=-\He_3(x)$ and $\He_b(-x)=(-1)^b\He_b(x)$, while
$q_a$ has parity $(-1)^{a-1}$; the integrand changes sign unless $a+b$ is odd.
\end{proof}

\begin{lemma}[Two--dimensional reduction]\label{lem:reduction}
Let
\begin{equation}\label{eq:rho}
   \rho(t)=\frac{t-s_3^2t^3+s_5^2t^5}{V}.
\end{equation}
Then
\begin{equation}\label{eq:Hred}
   H(t)=\frac{\pi}{2}\sum_{a,b\ge0}A_{a,b}^2\,\rho(t)^a\,(-t)^b .
\end{equation}
Consequently, for every $m\ge1$,
\begin{equation}\label{eq:bm}
   b_m=\frac{\pi}{2}\sum_{\substack{a,b\ge0\\a+b\ \mathrm{odd}}}
        (-1)^b\,A_{a,b}^2\;[t^{m-b}]\,\rho(t)^a .
\end{equation}
\end{lemma}

\begin{proof}
Expand $F$ in the orthonormal Hermite basis of $(W,X)$:
$F=\sum_{a,b}A_{a,b}\He_a(W)\He_b(X)$.  Writing $f$ and $g$ through the same $F$,
the variable $g$ uses $W'$ with the opposite degree--$3$ sign and the
degree--$5$ alignment, so the $W$--pair has correlation $\rho(t)$ (the linear
combination of $t,t^3,t^5$ with weights $1/V,-s_3^2/V,s_5^2/V$), while the
$X$--pair has correlation $t$ entering with the sign $(-1)^b$ from
$\He_b(-X)$.  The Hermite covariance identity gives
$\E[\He_a(W)\He_a(W')]=\rho(t)^a$ and
$\E[\He_b(X)\He_b(-X')]=(-1)^b t^b$, and \eqref{eq:Hred} follows.  Reading off
$[t^m]$ and using $\rho(t)=O(t)$ gives \eqref{eq:bm}.
\end{proof}

Each $A_{a,b}$ is a one--dimensional integral of an explicit entire function
(\cref{lem:qa}); \cref{sec:formulas} records how it is enclosed.  The squared
coefficients carry all the mass: by Parseval $\sum_{a,b}A_{a,b}^2=\E[F^2]=1$.

\begin{lemma}[Low coefficients]\label{lem:low}
With the enclosures of \cref{sec:formulas},
\[
   b_1\ge0.8815738220,\qquad
   b_3=3.4538733\cdot10^{-9},\qquad
   b_5=-9.0435536\cdot10^{-11},
\]
and the partial sums through degree $251$ satisfy
\begin{equation}\label{eq:heads}
   \sum_{\substack{3\le m\le251\\m\ \mathrm{odd}}}\abs{b_m}\le1.1328860\cdot10^{-5},
   \qquad
   \sum_{\substack{3\le m\le251\\m\ \mathrm{odd}}}m\abs{b_m}\le1.5737642\cdot10^{-4}.
\end{equation}
\end{lemma}

The near--vanishing of $b_3$ and $b_5$ is the design goal: series reversion gives
$A_3=-b_3/b_1^4$ and $A_5=(3b_3^2-b_1b_5)/b_1^7$, so killing $b_3,b_5$ kills the
first two nonlinear inverse obstructions.

\section{The coefficient tail through a Sobolev norm}\label{sec:sobolev}

Let $\Dop=t\,\frac{d}{dt}$ be the Euler operator, $\Dop\,t^m=m\,t^m$.  For the odd
series $H(t)=\sum_{m\ \mathrm{odd}}b_mt^m$ define
\begin{equation}\label{eq:Bs}
   B_s^2=\sum_{\substack{m\ge1\\m\ \mathrm{odd}}}m^{2s}\,b_m^2
        =\frac{1}{2\pi}\int_0^{2\pi}\bigl\lvert \Dop^sH(e^{i\theta})\bigr\rvert^2\,d\theta .
\end{equation}
The two expressions agree by Parseval on the circle $\T$.  Crucially, $B_s$ sees
the \emph{signed} coefficients $b_m$, not the positive Hermite mass; a triangle
inequality applied to \eqref{eq:Hred} would lose the cancellation that makes the
$b_m$ small.

\begin{lemma}[Sobolev tail]\label{lem:tail}
Let $N$ be odd.  Then
\[
   \sum_{m>N}\abs{b_m}\le B_3\Bigl(\sum_{\substack{m>N\\m\ \mathrm{odd}}}m^{-6}\Bigr)^{1/2},
   \qquad
   \sum_{m>N}m\abs{b_m}\le B_3\Bigl(\sum_{\substack{m>N\\m\ \mathrm{odd}}}m^{-4}\Bigr)^{1/2}.
\]
\end{lemma}

\begin{proof}
Write $\abs{b_m}=m^{-3}\cdot m^3\abs{b_m}$ and $m\abs{b_m}=m^{-2}\cdot m^3\abs{b_m}$
and apply Cauchy--Schwarz, using $\sum_{m>N}m^6b_m^2\le B_3^2$.
\end{proof}

The power--sum factors are tiny.  With $N=251$,
\begin{equation}\label{eq:tailfac}
   \Bigl(\sum_{\substack{m>251\\\mathrm{odd}}}m^{-6}\Bigr)^{1/2}\le3.13677\cdot10^{-7},
   \qquad
   \Bigl(\sum_{\substack{m>251\\\mathrm{odd}}}m^{-4}\Bigr)^{1/2}\le1.020511\cdot10^{-4};
\end{equation}
each odd power sum is enclosed by a finite head plus a decreasing--integral
remainder.  Thus even a crude bound on $B_3$ suffices, which is fortunate
because $B_3$ is the only hard quantity.

\section{Bounding the Sobolev norm}\label{sec:b3}

By \eqref{eq:Bs} it is enough to bound $\Dop^3H(e^{i\theta})$ on the circle.  We
realize $\Dop^3H$ as a trace of a two--dimensional differential kernel and apply
an averaged inequality.

\subsection{A combined kernel}

For $\abs r<1$ let $\Kr_r$ be the centered bivariate Gaussian density of
correlation $r$,
\begin{equation}\label{eq:Krkernel}
   \Kr_r(u,v)=\frac{1}{2\pi\sqrt{1-r^2}}\,
   \exp\!\Bigl(-\frac{u^2-2ruv+v^2}{2(1-r^2)}\Bigr),
\end{equation}
and set $h(x)=\He_3(x)$, so $\alpha h(x)$ is the threshold of
\cref{def:Aab}.  Define the base kernel
$\Psi_r(x,y)=\Kr_r(\alpha h(x),\alpha h(y))$, and, with $\lambda=-t$ and
$\Dop_t=t\,\partial_t$, the kernels
\begin{equation}\label{eq:phik}
   \Phi_0=\Psi_{\rho(t)},\qquad
   \Phi_{k+1}=\Dop_t\Phi_k+\lambda\,\partial_x\partial_y\Phi_k\quad(k\ge0).
\end{equation}
The combined kernel $\Phi_3$ is a single function of $(x,y)$ (with $t$ a
parameter), obtained by carrying out the recursion \eqref{eq:phik} before any
norm is taken.

\begin{lemma}[Trace representation]\label{lem:trace}
For $\abs t=1$,
\[
   \Dop^3H(t)=T_{-t}[\Phi_3(t;\,\cdot\,,\cdot\,)],
\]
where, for a kernel $\Psi$ and $\abs\lambda=1$, $T_\lambda$ is the bounded linear
functional satisfying
\begin{equation}\label{eq:traceineq}
   \abs{T_\lambda[\Psi]}\le\frac{\pi}{\sqrt6}
   \Bigl(\norm{\Psi}_{L^2(\mu)}^2
        +4\,\norm{\partial_x\partial_y\Psi}_{L^2(\mu)}^2\Bigr)^{1/2},
\end{equation}
with $\mu$ the standard Gaussian measure on $\R^2$.
\end{lemma}

\begin{proof}[Proof sketch]
Differentiating \eqref{eq:Hred} as $H(t)=\tfrac\pi2\,\Psi^{\sharp}(\rho(t),-t)$,
where $\Psi^{\sharp}(r,s)=\E[F(W,X)F(W',X')]$ at $W$--correlation $r$ and
$X$--correlation $s$, the Euler operator acts by the chain rule.  Price's theorem
(Gaussian integration by parts) identifies $\partial_s$ with $\partial_x\partial_y$
on the kernel and $\partial_r$ with the kernel's $r$--derivative; collecting the
terms produces exactly the recursion \eqref{eq:phik} with $\lambda=\Dop_t(-t)=-t$.
The remaining $W$--integration against the (degenerate, $\abs\lambda=1$) Gaussian
of correlation $-t$ is the functional $T_{-t}$.  Writing that integration as an
$L^2(\mu)$ pairing and using $\norm{h'}_{L^2}$ and the second--order Gaussian
weight gives \eqref{eq:traceineq}; the constant $\pi/\sqrt6$ comes from the
normalization of $\He_3$ in \eqref{eq:He3} together with the $\pi/2$ prefactor.
\end{proof}

The recursion \eqref{eq:phik} expands $\Phi_3$ as a finite sum over a
nine--term index set; the explicit prefactors and kernels are in
\cref{sec:formulas}.

\subsection{An averaged Minkowski bound}

From \cref{lem:trace} and \eqref{eq:Bs},
\begin{equation}\label{eq:B3avg}
   B_3^2=\frac{1}{2\pi}\int_0^{2\pi}\abs{\Dop^3H(e^{i\theta})}^2\,d\theta
   \le\frac{\pi^2}{6}\,\overline S,
   \qquad
   \overline S:=\frac{1}{2\pi}\int_0^{2\pi}S(\theta)\,d\theta,
\end{equation}
where $S(\theta)=\norm{\Phi_3}_{L^2(\mu)}^2+4\norm{\partial_x\partial_y\Phi_3}_{L^2(\mu)}^2$
and $\Phi_3$ is taken at $t=e^{i\theta}$.

\begin{lemma}[Averaged Minkowski]\label{lem:mink}
Write $\Phi_3=\sum_{p,a}c_{p,a}(t)\,G_{p,a}$ for the expansion produced by
\eqref{eq:phik}, with scalar prefactors $c_{p,a}$ and kernels $G_{p,a}$
(\cref{sec:formulas}).  For each pair set
\[
   Q^{(0)}_{p,a}=\frac{1}{2\pi}\!\int_0^{2\pi}\!\abs{c_{p,a}(e^{i\theta})}^2
            \norm{G_{p,a}}_{L^2(\mu)}^2 d\theta,
   \quad
   Q^{(1)}_{p,a}=\frac{1}{2\pi}\!\int_0^{2\pi}\!\abs{c_{p,a}(e^{i\theta})}^2
            \norm{\partial_x\partial_yG_{p,a}}_{L^2(\mu)}^2 d\theta .
\]
Then
\begin{equation}\label{eq:Smink}
   \overline S\le
   \Bigl(\sum_{p,a}\sqrt{Q^{(0)}_{p,a}}\Bigr)^{2}
   +4\Bigl(\sum_{p,a}\sqrt{Q^{(1)}_{p,a}}\Bigr)^{2}.
\end{equation}
\end{lemma}

\begin{proof}
By the triangle inequality in $L^2(\mu)$,
$\norm{\Phi_3}_{L^2(\mu)}\le\sum_{p,a}\abs{c_{p,a}}\norm{G_{p,a}}_{L^2(\mu)}$
pointwise in $\theta$, and likewise for $\partial_x\partial_y$.  Taking the
$L^2(d\theta/2\pi)$ norm of each sum and applying the triangle inequality there
gives $\overline{\norm{\Phi_3}^2}\le(\sum_{p,a}\sqrt{Q^{(0)}_{p,a}})^2$ and the
analogous bound for the second term; add them.
\end{proof}

Each $G_{p,a}$ is, on the real axis, a polynomial times a super--Gaussian, so the
$L^2(\mu)$ norms $\norm{G_{p,a}}^2$ and $\norm{\partial_x\partial_yG_{p,a}}^2$ are
honest two--dimensional Gaussian integrals; we enclose them rigorously and then
enclose the circle averages $Q^{(\cdot)}_{p,a}$ by a composite Gauss--Legendre
rule with a certified remainder (\cref{sec:formulas}).

\begin{lemma}[Certified Sobolev bound]\label{lem:B3val}
$\overline S\le1803.317$, and therefore
\[
   B_3\le\frac{\pi}{\sqrt6}\sqrt{1803.317}\le 54.4641 .
\]
\end{lemma}

The per--pair contributions, certified through the quadrature, are listed in
\cref{tab:Q}.  The bound $54.4641$ is generous: the true value of $B_3$ is about
$0.882$.  It does not need to be tight, because of the small tail factors
\eqref{eq:tailfac}.

\begin{table}[t]\centering
\small
\begin{tabular}{c r r}
\toprule
$(p,a)$ & $\sqrt{Q^{(0)}_{p,a}}$ & $\sqrt{Q^{(1)}_{p,a}}$ \\
\midrule
$(0,1)$ & $0.04286765$ & $0.10127946$ \\
$(0,2)$ & $0.30383835$ & $1.10543141$ \\
$(0,3)$ & $0.36847714$ & $7.40413211$ \\
$(1,0)$ & $2.67546361$ & $0.10772752$ \\
$(1,1)$ & $0.19214422$ & $0.85400602$ \\
$(1,2)$ & $0.51708607$ & $6.83760791$ \\
$(2,0)$ & $2.30597621$ & $0.29542681$ \\
$(2,1)$ & $0.41663640$ & $3.10773691$ \\
$(3,0)$ & $2.36356861$ & $0.91665956$ \\
\bottomrule
\end{tabular}
\caption{Certified upper bounds for the per--pair quantities of
\cref{lem:mink}; substituting into \eqref{eq:Smink} gives
$\overline S\le1803.317$.}
\label{tab:Q}
\end{table}

\section{Assembling the certificate}\label{sec:certificate}

We collect the pieces at $N=251$.  Combining \cref{lem:low,lem:tail} with
\eqref{eq:tailfac} and \cref{lem:B3val},
\[
   \eta_0\le\underbrace{1.1328860\cdot10^{-5}}_{\text{head}}
            +\underbrace{54.4641\cdot3.13677\cdot10^{-7}}_{\text{tail}}
      \le 2.84130\cdot10^{-5},
\]
\[
   \eta_1\le\underbrace{1.5737642\cdot10^{-4}}_{\text{head}}
            +\underbrace{54.4641\cdot1.020511\cdot10^{-4}}_{\text{tail}}
      \le 5.71549\cdot10^{-3}.
\]
With $b_1\ge0.8815738220496$ this gives $\eta_1<b_1$ and
\[
   \gam+\eta_0\le 0.881545409+2.84130\cdot10^{-5}
              = 0.8815738220 < b_1 .
\]
\cref{lem:wiener} applies with $\gam=0.881545409$, so
\[
   \KG\le\frac{\pi}{2\gam}\le 1.7818666069360661 .
\]
Since $\pi/(2\rhostar)=1.7822139781913691$, the improvement is at least
$3.4737\cdot10^{-4}$.  This proves \cref{thm:main}.\qed

\begin{remark}[Finite dimension]\label{rem:transfer}
\cref{def:H} is the fixed--degree limit of the finite scheme in which $R_3,R_5$
are replaced by normalized sums $N^{-1/2}\sum_{j\le N}\He_3(U_j)$ and
$N^{-1/2}\sum_{j\le N}\He_5(V_j)$ over independent Gaussians; the pair
$(f_N,g_N)$ then has dimension $d_N=2+2N$.  For each fixed $M$ the coefficients
$(b_{1,N},\dots,b_{M,N})$ converge to $(b_1,\dots,b_M)$ as $N\to\infty$, the
errors being collision terms of total $L^2$ mass $O_M(N^{-1})$.  Because
\cref{thm:main} uses finitely many coefficient inequalities with a positive
margin, the same inequalities hold for $(f_N,g_N)$ for all large $N$; the
limiting bound is therefore attained, up to the margin, by a genuine
finite--dimensional Krivine scheme.
\end{remark}

\section{Computational formulas}\label{sec:formulas}

This section lists every explicit object used above, in the form in which it is
enclosed by interval arithmetic.  All real numbers below are decimal inputs and
all derived quantities are outward--rounded \texttt{Arb} balls.

\paragraph{Parameters.}
$V=1+s_3^2+s_5^2$, $\alpha=\eta/\sqrt V$, and
$\rho(t)=(t-s_3^2t^3+s_5^2t^5)/V$ as in \eqref{eq:VW},\eqref{eq:rho}.

\paragraph{Coefficient integrals.}
$A_{a,b}=\int_\R \He_b(x)\,q_a(\alpha\,\He_3(x))\,\varphi(x)\,dx$, with $q_a$ from
\cref{lem:qa}.  The integrand is entire; on the real axis it is a polynomial
times $e^{-x^2/2}$ and decays super--Gaussianly, so it is enclosed by Arb's
adaptive integrator on a finite interval $[-L,L]$ plus the Gaussian tail
$\int_{\abs x>L}\abs{\He_b}\,\varphi\le$ a closed bound (using $\abs{q_a}\le1$).
Then $b_m$ is the interval sum \eqref{eq:bm}, with $[t^{m-b}]\rho(t)^a$ a finite
convolution of $(1,-s_3^2,s_5^2)/V$.  Absolute values in $\eta_0,\eta_1$ are taken
on enclosures bounded away from $0$.

\paragraph{Combined kernel.}
The recursion \eqref{eq:phik} produces $\Phi_3=\sum_{(p,a)}c_{p,a}(t)\,G_{p,a}$
over the nine pairs of \cref{tab:Q}.  The scalar prefactors $c_{p,a}(t)$ are
polynomials in $t$, generated by $c_{0,0}=1$ and, three times, the substitutions
\[
   c_{p,a}\ \longmapsto\
   \Dop_tc_{p,a}\ \text{ (same $(p,a)$)},\quad
   (\Dop_t\rho)\,c_{p,a}\ \text{ at }(p+1,a),\quad
   (-t)\,c_{p,a}\ \text{ at }(p,a+1).
\]
The kernels are, with $\varphi_r=\Kr_r$ from \eqref{eq:Krkernel} evaluated at
$u=\alpha h(x),\ v=\alpha h(y)$ and $r=\rho(t)$,
\[
   G_{p,0}=2\pi\,\partial_r^{\,p-1}\varphi_r\ (p\ge1),\qquad
   G_{p,a}=\partial_x^{a-1}\partial_y^{a-1}
           \bigl(2\pi\alpha^2\,h'(x)h'(y)\,\partial_r^{\,p}\varphi_r\bigr)\ (a\ge1).
\]
Each $\partial_r^{\,p}\partial_x^i\partial_y^j\varphi_r$ equals a polynomial in
$(r,u,v,1-r^2)$ times $\varphi_r$, with integer coefficients computed by a fixed
operator recurrence; this keeps the kernels exact.

\paragraph{Gaussian norms.}
Writing $u=\alpha h(x)$, the squared modulus of a kernel against $\mu$ is
\[
   \norm{G_{p,a}}_{L^2(\mu)}^2=\iint_{\R^2}\abs{G_{p,a}(x,y)}^2\,
        \frac{e^{-(x^2+y^2)/2}}{2\pi}\,dx\,dy,
\]
a positive integral of a polynomial times $e^{-x^2/2-y^2/2}$ and the bounded factor
\[
   \exp\!\bigl(-A\alpha^2(h(x)^2+h(y)^2)+2B\alpha^2h(x)h(y)\bigr),
   \qquad A=\Re(1/d),\ \ B=\Re(r/d),\ \ d=1-r^2 .
\]
Using $h(x)^3=\sqrt6\,h(x)+3x$, every mode reduces to a combination of the
one--dimensional base moments
$I_{\varepsilon,q}(A)=\E[X^\varepsilon h(X)^q e^{-A\alpha^2h(X)^2}]$,
$\varepsilon\in\{0,1,2\}$, each enclosed by Arb's integrator; the cross term is
expanded in a convergent series with a Chernoff tail.

\paragraph{Circle quadrature and its remainder.}
The averages $Q^{(\cdot)}_{p,a}$ are computed by a composite four--point
Gauss--Legendre rule on $[0,2\pi]$ with $M$ panels of the certified point values
$\abs{c_{p,a}}^2\norm{G_{p,a}}^2$.  Writing $g(\theta)$ for the integrand, the
rule's error is at most $C_4(2\pi)^8M^{-9}\sum_j D_{8,j}$ with
$C_4=(4!)^4/(9\,(8!)^3)$, where each panel bound $D_{8,j}\ge\sup_{I_j}\abs{g^{(8)}}$
is obtained \emph{without} cancellation from
\[
   \abs{g^{(8)}}\le\sum_{a=0}^{8}\binom8a\sqrt{V_aV_{8-a}},\qquad
   V_a(\theta)=\norm{\partial_\theta^a(c_{p,a}\,G_{p,a})}_{L^2(\mu)}^2 ,
\]
and $V_a\le\abs d^{-1}\bigl(\sum_\mu\abs{p_{a,\mu}(\theta)}\sqrt{L_\mu L_{\mu'}}\bigr)^2$
by the elementary inequality $2BUV\le(A-\abs B)(U^2+V^2)$, valid because
$A-\abs B=\Re(1/d)-\abs{\Re(r/d)}\ge\tfrac12$ on the unit circle (this is verified
on all panels).  Here $L_\mu$ are precomputed positive one--dimensional integrals
and $\abs{p_{a,\mu}}$ are bounded by evaluating the $\theta$--derivative
coefficients on each panel as an interval.  At $M=2048$ the total remainder is at
most $9.1\cdot10^{-4}$ per pair, negligible against $\overline S$.

\section{Implementation notes}\label{sec:impl}

The following are not needed to state the proof but record why the certificate is
organized as it is; the corresponding code lives in the repository.

\begin{enumerate}[leftmargin=1.6em]
\item \textbf{Why a point evaluation, not interval boxes.}  The integrand of a
Gaussian norm decays like $e^{-c\,x^6}$ on the real axis but \emph{grows} like
$e^{+c\abs z^6}$ off it, because $h(z)^2\sim z^6/6$ has negative real part in
sectors.  Interval boxes in $(x,y)$, or a complex contour, see this far--field
growth and blow up.  Every integral is therefore evaluated at a real point;
\texttt{Arb}'s adaptive integrator works because its per--subinterval error model
only probes small ellipses that never reach the far field.

\item \textbf{Why the Sobolev tail, not an inverse tail.}  Bounding $\sum\abs{b_m}$
through $B_3=\norm{\Dop^3H}_{L^2}$ keeps the cancellation: $B_3$ is built from the
already--small signed $b_m$.  A direct Hermite bound, or a complex Cauchy estimate
of $H^{-1}$ near $t=1$, overshoots by many orders of magnitude.

\item \textbf{Why the combined kernel.}  Forming $\Phi_3$ by the recursion
\eqref{eq:phik} \emph{before} taking norms is essential: expanding
$\Dop^3H$ as a sum of separately bounded terms discards a nine--fold cancellation
and inflates $\overline S$ from $\approx1.8\cdot10^3$ to the order of $10^2$ per
term.

\item \textbf{Why the cancellation--free $D_8$.}  The exact eighth $\theta$--derivative
$g^{(8)}$ is a signed combination of inner products whose terms are $\sim10^{40}$
and cancel to $\sim10^{10}$; reproducing this in interval arithmetic needs
$\sim140$ digits near $\theta=0,\pi$.  The Minkowski bound of
\cref{sec:formulas} avoids the signed sum entirely: every $V_a$ is a positive
quantity, immune to the cancellation, at the cost of a looser $D_8$ that the
$M^{-9}$ factor absorbs.

\item \textbf{Precision.}  The base moments overflow a $60$--digit mantissa for
$q\gtrsim230$; the circle nodes near $\theta=0,\pi$ use a longer Chernoff
truncation and $200$--bit precision.  The reported headline numbers are stable
under further refinement.
\end{enumerate}

\end{document}
